\chapter{Sprint 2: Analytics and Uncertainty}
\chaptertoc

\section{Introduction}
Sprint 2 transformed perception outputs into operational queue indicators. The focus moved from ``who is detected'' to ``how the queue behaves'' through arrival/service-rate estimation, wait-time computation, and stability-aware decision logic.

\section{Sprint Backlog}
Sprint 2 backlog priorities were:
\begin{itemize}
    \item Implement event-based queue rate estimation from tracked IDs in zone.
    \item Compute wait-time estimates with explicit stable/unstable queue handling.
    \item Add denominator-safety guards and fallback wait-time logic.
    \item Expose a consistent runtime metric contract for dashboard and webhook consumption.
\end{itemize}

\section{Design}

\subsection{Component Diagram}
\begin{figure}[ht]
    \centering
    \fbox{\parbox[c][6cm][c]{0.85\textwidth}{\centering Placeholder for Sprint 2 component diagram}}
    \caption{Sprint 2 components for queue analysis and stability handling.}
    \label{fig:sprint2_component}
\end{figure}

At component level, \texttt{QueueAnalyzer} receives in-zone tracker IDs and produces \texttt{QueueMetrics} consumed by logs, webhook payloads, and dashboard updates.

\subsection{Data Flow Diagram}
\begin{figure}[ht]
    \centering
    \fbox{\parbox[c][6cm][c]{0.85\textwidth}{\centering Placeholder for Sprint 2 data flow diagram}}
    \caption{Data flow from tracked events to queue metrics and stability state.}
    \label{fig:sprint2_data_flow}
\end{figure}

Data flow is event-centric: per-frame ID transitions produce arrival/departure events, event windows produce rates $(\lambda,\mu)$, and rates with occupancy produce wait-time estimates plus stable/unstable state labels.

\subsection{Sequence Diagrams}
The implemented sequence is:
\texttt{QueueAnalyzer.update()} computes counts and rates, evaluates queue stability, and applies either the stable wait-time formula or fallback wait-time branch before publishing \texttt{QueueMetrics}.

\section{Implementation}
\subsection{Queue Event Windows and Rate Estimation}
In \texttt{backend/src/queue\_analyzer.py}, arrivals and departures are tracked as timestamp deques over sliding windows. Let $N_a$ and $N_d$ be events in windows $T_a$ and $T_d$:
\[
\lambda = \frac{N_a}{T_a}, \qquad \mu = \frac{N_d}{T_d}
\]
The implementation enforces a minimum event count before returning non-zero rates, reducing unstable early-window estimates.

\subsection{Wait-Time Logic and Stability Handling}
The runtime marks queue stability with:
\[
\texttt{queue\_stable} = (\lambda < \mu) \land (\mu > 0)
\]
The wait-time logic follows two branches:
\begin{itemize}
    \item Stable queue: $W = \frac{1}{\mu-\lambda}$
    \item Unstable but service observed: fallback $W = \frac{L}{\mu}$, where $L$ is people in zone
\end{itemize}
When no service events are observed, the implementation currently returns a zero fallback.

\subsection{Runtime Contract for Monitoring and Automation}
The final \texttt{QueueMetrics} object includes:
\begin{itemize}
    \item point estimates ($\lambda$, $\mu$, $W$),
    \item people-in-zone occupancy,
    \item stability flag for safe downstream interpretation.
\end{itemize}
These fields are propagated to logs, overlays, webhook payloads, and dashboard events.

\section{Conclusion}
Sprint 2 delivered the analytical intelligence of the project: event-to-rate modeling, stable/unstable wait-time semantics, and a consistent runtime metric contract. This provided the operational core required for API exposure and external automation in the next sprint.
